\section{Metoda}

\subsection{Oznaczenia}
Złożenia czytamy od prawej (prawy czynnik działa pierwszy). Niech
\[
t_f(x,y)=(x,\,y+f(x)),\qquad u_g(x,y)=(x+g(y),\,y),\qquad \sigma(x,y)=(y,x),
\]
gdzie $f,g\in\Q[t]$. Przez $\mathrm{Aff}$ oznaczamy grupę odwracalnych odwzorowań afinicznych
nad $\Q$, przez $J$ grupę odwzorowań trójkątnych $(\alpha x+c,\ \beta y+h(x))$, $\alpha\beta\neq0$,
$h\in\Q[x]$, a $B=\mathrm{Aff}\cap J$. Zauważmy, że $\sigma t_g\sigma=u_g$.

\subsection{Struktura grupy automorfizmów}
Podstawą jest klasyczne twierdzenie, którego \emph{nie dowodzimy}, lecz cytujemy.

\begin{twierdzenie}[Jung 1942, van der Kulk 1953; cytowane]\label{tw:jvdk}
Dla dowolnego ciała $k$ grupa $\Aut(k^2)$ jest iloczynem amalgamowanym
$\mathrm{Aff}*_B J$. W szczególności każdy nieafiniczny automorfizm $F$ ma słowo zredukowane
$F=L_kT_k L_{k-1}\cdots T_1L_0$ z $T_i\in J\setminus B$, $L_i\in\mathrm{Aff}$,
$L_1,\dots,L_{k-1}\notin B$, a długość $k$ jest niezmiennikiem $F$.
\end{twierdzenie}

\begin{lemat}[postać standardowa; UDOWODNIONE, iter01/zad3 i iter02/zad2]\label{lem:postac}
Każdy nieafiniczny $F\in\Aut(\Q^2)$ da się zapisać jako
\[
F=A\circ\sigma t_{f_k}\sigma t_{f_{k-1}}\cdots\sigma t_{f_1}\circ L,\qquad A,L\in\mathrm{Aff},\ f_i\in\Q[x],\ \deg f_i\ge2,
\]
i wtedy $\deg F=\prod_{i=1}^k\deg f_i$.
\end{lemat}
\begin{proof}[Szkic]
Element $J\setminus B$ ma postać $d\circ t_h$ z $d$ diagonalno-afinicznym i $\deg h\ge2$.
Rozkład Bruhata $\mathrm{Aff}=B\cup B\sigma B$ (używa tylko działań ciała, więc działa nad $\Q$)
pozwala zapisać $L_i=b_i\sigma b_i'$; bloki $b_i'T_ib_{i-1}$ leżą w $J\setminus B$, a ich części
diagonalne przesuwamy w lewo przez $\sigma$, bo $\sigma d\sigma\in B$ dla diagonalnego $d$
(dla dowolnego elementu $B$ to nie zachodzi -- doprecyzowanie z recenzji iter02/zad2).
Dla stopnia: jeśli $G_{j-1}=(P,R)$ z $\deg P=D_{j-1}$, $\deg R\le D_{j-1}$, to
$G_j=\sigma t_{f_j}G_{j-1}=(R+f_j(P),P)$, a forma wiodąca $f_j(P)$ ma stopień
$\deg f_j\cdot D_{j-1}>D_{j-1}\ge\deg R$, więc nie może się skrócić. Indukcja daje
$\deg F=\prod\deg f_i$; złożenie z afinicznym $A$ stopnia nie zmienia.
\end{proof}

\subsection{Warunki interpolacyjne dla słów długości 2}
Dla $k=2$ mamy $F=A\circ u_g\circ t_f\circ L$. Piszemy $L(p)=(a(p),b(p))$,
$A^{-1}(q)=(c(q),e(q))$, gdzie $(a,b)$ i $(c,e)$ są parami liniowo niezależnych funkcjonałów
afinicznych. Ponieważ $u_gt_f(a,b)=(a+g(b+f(a)),\,b+f(a))$, warunki $F(p_i)=q_i$ są
\emph{równoważne} układowi
\begin{equation}\label{eq:warunki}
e(q_i)-b(p_i)=f\bigl(a(p_i)\bigr),\qquad c(q_i)-a(p_i)=g\bigl(e(q_i)\bigr),\qquad i=1,\dots,5.
\end{equation}
Przy \emph{ustalonych} funkcjonałach $a$, $e$ układ \eqref{eq:warunki} jest liniowy w
pozostałych niewiadomych (współczynniki $b$, $c$, $f$, $g$). Wyrazy stałe $a$, $e$ i wyrazy liniowe
$f,g$ można wchłonąć w $b$, $c$ (normalizacje opisane w notatkach iter02/zad1 i iter02/zad2), więc
bez straty ogólności $a=\alpha_1x+\alpha_2y$, $e=\varepsilon_1x+\varepsilon_2y$.

\subsection{Metody obliczeniowe}
Wszystkie rachunki prowadzono dokładnie nad $\Q$ (Sage/Singular oraz sympy).
\begin{enumerate}
\item \textbf{Brak zgodnego funkcjonału.} Jądro macierzy $5\times6$ o wierszach
$[p_x,p_y,1,-q_x,-q_y,-1]$ (skrypty \texttt{iter01/zad3/skrypty/brak\_funkcjonalu.py},
\texttt{iter02/zad3/skrypty/brak\_funkcjonalu\_niezal.py}, \texttt{iter02/zad2/skrypty/kontrola.sage}).
\item \textbf{Konstrukcje stopnia 6.} Dla multistopni $(2,3)$ i $(3,2)$: przeglądanie małych
kierunków $(\alpha_1,\alpha_2)$, $(\varepsilon_1,\varepsilon_2)$ i rozwiązywanie liniowego układu
\eqref{eq:warunki} nad $\Q$ (\texttt{iter01/zad2/skrypty/szukaj.sage},
\texttt{iter02/zad3/skrypty/szukaj\_ladne.py}: wszystkie prymitywne kierunki całkowite o
współrzędnych $|\cdot|\le3$).
\item \textbf{Wykluczenie multistopnia $(2,2)$, metoda I} (iter01/zad1, \texttt{st4.sage};
powtórzona z jawnym przeglądem kart w iter02/zad2, \texttt{kontrola.sage}). Przy ustalonych
$a,e$ istnienie rozwiązania pierwszej grupy równań \eqref{eq:warunki} z $\deg f\le2$ wymaga
$\det M_1=0$, $M_1=[p_x,p_y,1,a(p)^2\mid e(q)]$ (wiersze $i=1..5$), a drugiej -- $\det M_2=0$,
$M_2=[q_x,q_y,1,e(q)^2\mid a(p)]$. Oba wielomiany są dwujednorodne w $(\alpha,\varepsilon)$;
ideał nasycono względem $\alpha=0$ i $\varepsilon=0$ i zbadano wszystkie karty afiniczne
$\mathbb P^1\times\mathbb P^1$.
\item \textbf{Wykluczenie multistopnia $(2,2)$, metoda II} (iter02/zad1, \texttt{gb22.sage}).
Pełny ideał wszystkich współczynników po normalizacji $f=x^2$, $g=g_2y^2$, warunki niezerowości
przez zmienne pomocnicze ($u\det(a,b)=1$ itd.), baza Gröbnera i eliminacja do niezmiennika
$s=\alpha_1^3g_2$. \emph{Test kontrolny}: po ustaleniu $g_2=1$ zbudowano jawne rozwiązania nad
ciałami liczbowymi i sprawdzono symbolicznie, że dają automorfizmy stopnia 4 realizujące bijekcję
-- co pokazuje, że metoda nie jest ,,ślepa''.
\end{enumerate}
Rozwiązania sprawdzał niezależny weryfikator (\texttt{narzedzia/verify.py}), który dokładnie
sprawdza $F(p_i)=q_i$, istnienie wymiernej odwrotności wielomianowej i stopień.
